\documentclass[11pt]{article}
\usepackage{mathblatt}
\begin{document}
\weit
\typeout{PROBE-BEGIN rekonstruktion-von-bestaenden-e1-k3-s3-v2}
\begin{aufgabe}{\detokenize{rekonstruktion-von-bestaenden-e1-k3-s3-v2}}
\begin{teile}
\teil Probe

\begin{ksys}[xmin=0,xmax=6,ymin=0,ymax=4,xlabel=t,ylabel=r] \funktion{1+0.5*\x}{r} \flaeche{1+0.5*\x}{2}{5}{A} \end{ksys}
\end{teile}
\end{aufgabe}
\typeout{PROBE-END rekonstruktion-von-bestaenden-e1-k3-s3-v2}
\typeout{PROBE-BEGIN uneigentliche-integrale-e2-k1-s0-v1}
\begin{aufgabe}{\detokenize{uneigentliche-integrale-e2-k1-s0-v1}}
\begin{teile}
\teil Probe

\begin{ksys}[xmin=0,xmax=10,ymin=0,ymax=3]\funktion{3*exp(-\x)}{f}\flaeche{3*exp(-\x)}{5}{10}{R}\end{ksys}
\end{teile}
\end{aufgabe}
\typeout{PROBE-END uneigentliche-integrale-e2-k1-s0-v1}
\typeout{PROBE-BEGIN uneigentliche-integrale-e2-k1-s0-v2}
\begin{aufgabe}{\detokenize{uneigentliche-integrale-e2-k1-s0-v2}}
\begin{teile}
\teil Probe

\begin{ksys}[xmin=0,xmax=10,ymin=0,ymax=3]\funktion{2*exp(-0.2*\x)}{f}\flaeche{2*exp(-0.2*\x)}{3}{10}{R}\end{ksys}
\end{teile}
\end{aufgabe}
\typeout{PROBE-END uneigentliche-integrale-e2-k1-s0-v2}
\typeout{PROBE-BEGIN uneigentliche-integrale-e2-k1-s0-v3}
\begin{aufgabe}{\detokenize{uneigentliche-integrale-e2-k1-s0-v3}}
\begin{teile}
\teil Probe

\begin{ksys}[xmin=0,xmax=6,ymin=0,ymax=4]\funktion{4*exp(-2*\x)}{f}\flaeche{4*exp(-2*\x)}{3}{6}{R}\end{ksys}
\end{teile}
\end{aufgabe}
\typeout{PROBE-END uneigentliche-integrale-e2-k1-s0-v3}
\typeout{PROBE-BEGIN uneigentliche-integrale-e2-k1-s0-v4}
\begin{aufgabe}{\detokenize{uneigentliche-integrale-e2-k1-s0-v4}}
\begin{teile}
\teil Probe

\begin{ksys}[xmin=0,xmax=10,ymin=0,ymax=3]\funktion{2/(\x+1)}{f}\flaeche{2/(\x+1)}{4}{10}{R}\end{ksys}
\end{teile}
\end{aufgabe}
\typeout{PROBE-END uneigentliche-integrale-e2-k1-s0-v4}
\end{document}
